We prove both identities by a direct factorization and finite linearity of coefficient extraction.

Work formally in a Laurent-series setting, for instance over the coefficient field \(K=\mathbb{Q}(\lambda_1,\lambda_2)\). The residue map
\[
\operatorname{Res}_{h=0}:K((h))\to K,\qquad \operatorname{Res}_{h=0}[f(h)]=[h^{-1}]f(h),
\]
is \(K\)-linear. In particular, for any finite set \(S\) and any Laurent series \(f_k(h)\),
\[
\sum_{k\in S}\operatorname{Res}_{h=0}[f_k(h)]
=
\operatorname{Res}_{h=0}\!\left[\sum_{k\in S} f_k(h)\right].
\]
The finiteness of \(S\) is the only point needed here; no convergence issue is involved.

\medskip\noindent\rule{\linewidth}{0.4pt}

\subsubsection*{The \(A\)-sum}

Let
\[
S_A=\{k\in\mathbb{Z}:1\le k\le l_1,\ k\equiv l_1\pmod 2\}.
\]
For \(k\in S_A\), compute \(P_A(h)q_A(h)^k\). By definition,
\[
q_A(h)^k
=
\frac{(h+\lambda_2)^k(h+\lambda_1-\lambda_2)^k(h-2\lambda_1)^k}
{h^k(h-\lambda_1+\lambda_2)^k(h-2\lambda_1-\lambda_2)^k}.
\]
Therefore
\begin{align*}
P_A(h)q_A(h)^k
&=
(h-\lambda_1)^2
(h+\lambda_2)^{l_2}
(h+\lambda_1-\lambda_2)^{l_3}
(h-\lambda_1+\lambda_2)^{l_1-l_2} \\
&\quad\times
(h-2\lambda_1-\lambda_2)^{l_1-l_3}
(h-2\lambda_1)^{-2-2l_1} \\
&\quad\times
\frac{(h+\lambda_2)^k(h+\lambda_1-\lambda_2)^k(h-2\lambda_1)^k}
{h^k(h-\lambda_1+\lambda_2)^k(h-2\lambda_1-\lambda_2)^k}.
\end{align*}
Combining like bases using the usual exponent laws for integer powers gives
\begin{align*}
P_A(h)q_A(h)^k
&=
h^{-k}(h-\lambda_1)^2
(h+\lambda_2)^{k+l_2}
(h+\lambda_1-\lambda_2)^{k+l_3} \\
&\quad\times
(h-\lambda_1+\lambda_2)^{l_1-l_2-k}
(h-2\lambda_1-\lambda_2)^{l_1-l_3-k}
(h-2\lambda_1)^{k-2-2l_1}.
\end{align*}
The right-hand side is exactly the integrand appearing in the definition of \(A(l_1,l_2,l_3,k)\). Hence
\[
A(l_1,l_2,l_3,k)
=
\operatorname{Res}_{h=0}\!\left[P_A(h)q_A(h)^k\right].
\]

Now sum over the parity-restricted finite set \(S_A\). By linearity of \(\operatorname{Res}_{h=0}\),
\begin{align*}
\sum_{\substack{1\le k\le l_1\\ k\equiv l_1\pmod 2}}
A(l_1,l_2,l_3,k)
&=
\sum_{k\in S_A}
\operatorname{Res}_{h=0}\!\left[P_A(h)q_A(h)^k\right] \\
&=
\operatorname{Res}_{h=0}\!\left[
\sum_{k\in S_A} P_A(h)q_A(h)^k
\right] \\
&=
\operatorname{Res}_{h=0}\!\left[
P_A(h)\sum_{k\in S_A}q_A(h)^k
\right].
\end{align*}
But by definition,
\[
\Phi_{l_1}(z)
=
\sum_{\substack{1\le k\le l_1\\ k\equiv l_1\pmod 2}} z^k,
\]
so substituting \(z=q_A(h)\) gives
\[
\sum_{k\in S_A}q_A(h)^k
=
\Phi_{l_1}(q_A(h)).
\]
Therefore
\[
\boxed{
\sum_{\substack{1\le k\le l_1\\ k\equiv l_1\pmod2}} A(l_1,l_2,l_3,k)
=
\operatorname{Res}_{h=0}\!\left[P_A(h)\,\Phi_{l_1}(q_A(h))\right].
}
\]

\medskip\noindent\rule{\linewidth}{0.4pt}

\subsubsection*{The \(B\)-sum}

Assume now that \(l_2\ge 1\), and define
\[
S_B=\{k\in\mathbb{Z}:1\le k\le l_2,\ k\equiv l_2\pmod 2\}.
\]
For \(k\in S_B\), we similarly compute
\[
q_B(h)^k
=
\frac{(h+\lambda_1)^k(h+\lambda_2-\lambda_1)^k(h-2\lambda_2)^k}
{h^k(h-\lambda_2+\lambda_1)^k(h-2\lambda_2-\lambda_1)^k}.
\]
Thus
\begin{align*}
P_B(h)q_B(h)^k
&=
(h-\lambda_2)^2
(h+\lambda_1)^{l_1}
(h+\lambda_2-\lambda_1)^{l_3}
(h-\lambda_2+\lambda_1)^{l_2-l_1} \\
&\quad\times
(h-2\lambda_2-\lambda_1)^{l_2-l_3}
(h-2\lambda_2)^{-2-2l_2} \\
&\quad\times
\frac{(h+\lambda_1)^k(h+\lambda_2-\lambda_1)^k(h-2\lambda_2)^k}
{h^k(h-\lambda_2+\lambda_1)^k(h-2\lambda_2-\lambda_1)^k}.
\end{align*}
Combining powers gives
\begin{align*}
P_B(h)q_B(h)^k
&=
h^{-k}(h-\lambda_2)^2
(h+\lambda_1)^{k+l_1}
(h+\lambda_2-\lambda_1)^{k+l_3} \\
&\quad\times
(h-\lambda_2+\lambda_1)^{l_2-l_1-k}
(h-2\lambda_2-\lambda_1)^{l_2-l_3-k}
(h-2\lambda_2)^{k-2-2l_2}.
\end{align*}
This is exactly the integrand defining \(B(l_1,l_2,l_3,k)\). Hence
\[
B(l_1,l_2,l_3,k)
=
\operatorname{Res}_{h=0}\!\left[P_B(h)q_B(h)^k\right].
\]

Using finite linearity of the residue,
\begin{align*}
\sum_{\substack{1\le k\le l_2\\ k\equiv l_2\pmod 2}}
B(l_1,l_2,l_3,k)
&=
\sum_{k\in S_B}
\operatorname{Res}_{h=0}\!\left[P_B(h)q_B(h)^k\right] \\
&=
\operatorname{Res}_{h=0}\!\left[
\sum_{k\in S_B}P_B(h)q_B(h)^k
\right] \\
&=
\operatorname{Res}_{h=0}\!\left[
P_B(h)\sum_{k\in S_B}q_B(h)^k
\right] \\
&=
\operatorname{Res}_{h=0}\!\left[
P_B(h)\Phi_{l_2}(q_B(h))
\right].
\end{align*}
Therefore
\[
\boxed{
\sum_{\substack{1\le k\le l_2\\ k\equiv l_2\pmod2}} B(l_1,l_2,l_3,k)
=
\operatorname{Res}_{h=0}\!\left[P_B(h)\,\Phi_{l_2}(q_B(h))\right].
}
\]

Both desired compressed residue identities follow.